\documentclass[11pt]{article}
\usepackage{amsmath, amssymb, amsthm}
\usepackage{geometry}
\geometry{margin=1in}

\title{Mathematical Spine for the Geometric Hallucination Taxonomy}
\author{Borealis S. Hedling}
\date{2026}

\begin{document}
\maketitle

\section*{Overview}
This document provides the formal mathematical backbone for the 
\textit{Geometric Hallucination Taxonomy (GHT)}. 
It defines the geometric operators used to classify hallucinations in 
structural AI systems, including curvature, spectral derivatives, holonomy, 
topological rank instability, and fusion-field nonlinearity.

Each operator is defined independently of model architecture, enabling 
future formal verification and Lean-based proof development.

\section{Curvature-Driven Hallucinations (CDH)}
Hallucinations arising from instability in the latent manifold's metric tensor 
$g_{ij}$ are modeled using the Riemann curvature tensor.

\subsection{Riemann Curvature Tensor}
Given a smooth manifold $(M, g)$ with Levi-Civita connection $\nabla$, the 
Riemann curvature tensor is defined as:
\[
R(X,Y)Z = \nabla_X \nabla_Y Z - \nabla_Y \nabla_X Z - \nabla_{[X,Y]} Z.
\]

In coordinates:
\[
R^i_{\ jkl} = \partial_k \Gamma^i_{\ jl} - \partial_l \Gamma^i_{\ jk}
+ \Gamma^i_{\ km}\Gamma^m_{\ jl} - \Gamma^i_{\ lm}\Gamma^m_{\ jk}.
\]

\subsection{Sectional Curvature}
Hallucination curvature signature:
\[
K(v) = R_{ijkl} v^i v^j v^k v^l.
\]

Positive blowout, negative collapse, and oscillatory curvature correspond to 
distinct hallucination modes.

\section{Spectral Cascade Hallucinations (SCH)}
Spectral instabilities are modeled using eigenvalue derivatives of the 
representation operator.

\subsection{Spectral Drift}
Let $\lambda_i(t)$ be the $i$-th eigenvalue over a representation trajectory.
Define:
\[
\Delta_\lambda = \frac{d\lambda_i}{dt}.
\]

Large $\Delta_\lambda$ magnitudes correspond to cascade onset.

\subsection{Spectral Collapse}
Rank loss is detected via:
\[
\text{rank}(A(t)) < \text{rank}(A(t-\epsilon)).
\]

\section{Holonomy Drift Hallucinations (HDH)}
Holonomy captures semantic drift arising from non-trivial parallel transport.

\subsection{Parallel Transport}
Given a loop $\gamma : [0,1] \to M$, parallel transport yields:
\[
P_\gamma : T_{\gamma(0)}M \to T_{\gamma(0)}M.
\]

\subsection{Holonomy Operator}
Hallucination drift signature:
\[
H = P_\gamma - I.
\]

Non-zero holonomy indicates rotational semantic drift.

\section{Topological Instability Hallucinations (TIH)}
Topological collapse is modeled using rank instability of the metric tensor.

\subsection{Metric Rank Operator}
Define:
\[
\tau = \text{rank}(g + \epsilon I).
\]

Rank drops correspond to collapse surfaces and discontinuities.

\section{Fusion-Field Hallucinations (FFH)}
Fusion-field hallucinations arise from nonlinear blending of latent concepts.

\subsection{Fusion Operator}
Let $x, y$ be latent representations. Define:
\[
F(x,y) = \alpha x + \beta y + \gamma \phi(x,y),
\]
where $\phi$ is a nonlinear interaction term.

Nonlinear fusion spikes correspond to chimera hallucinations.

\section{Future Lean Formalization}
The following definitions will be ported into Lean:

\begin{itemize}
    \item Riemann curvature tensor
    \item Sectional curvature operator
    \item Spectral derivative operator
    \item Holonomy operator
    \item Metric rank operator
    \item Fusion-field operator
\end{itemize}

These definitions will support theorem development for hallucination 
classification and operator stability.

\end{document}
